

\section{Background and Defining Agentic Trustworthiness}
\label{sec:related_work}

\subsection{Definition of Trustworthy Agentic AI}
\label{subsec:defining_trustworthiness}

Existing surveys treat trustworthiness as a vague umbrella or a long, unprioritised list~\cite{liu2024trustllm,wang2024trustworthyagent}. Building on emerging AI risk frameworks~\cite{nist2023airmf,eu2024aiact}, we adopt a deployment-oriented definition with five \emph{measurable} properties that cover the failure surface of an autonomous agent acting through tools in an open environment:

\begin{itemize}[leftmargin=*]
\item \textbf{(P1) Reliability.} Correct, complete outcomes on legitimate goals; failures include goal drift, hallucinated tool calls, and \emph{improper refusal} of benign requests (\S\ref{subsec:vulnerability_attribution}).
\item \textbf{(P2) Robustness.} Intended behaviour under adversarial inputs (prompt injection in tool outputs, retrieved documents, or user turns)~\cite{greshake2023indirect,debenedetti2024agentdojo} and under noisy or partial signals.
\item \textbf{(P3) Safety.} Respect for permission boundaries: avoiding unauthorised or irreversible actions and gating high-risk operations (writes, dispatch, payments) behind confirmation.
\item \textbf{(P4) Social-Ethical Alignment.} Fair, non-manipulative behaviour under social pressure---resisting coercion, role-play exploitation, discriminatory requests---and protecting third-party information~\cite{zhou2023sotopia,pan2023machiavelli}.
\item \textbf{(P5) Operational Integrity.} Graceful degradation under resource constraints, recovery from tool failures, adherence to step/budget envelopes, and audit-traceable trajectories.
\end{itemize}

We treat P1--P5 as the \emph{evaluation targets} of HAAF: trustworthiness is multi-dimensional, distributional, and deployment-conditioned---a property of the deployed system (weights + prompts + tools + guardrails), not the model alone~\cite{ganguli2022redteaming}. This definition governs the rest of the paper: \S\ref{subsec:overview}--\ref{subsec:layer4} map each layer to a subset of P1--P5, the taxonomy in \S\ref{subsec:vulnerability_attribution} assigns failures to properties, and the profile in \S\ref{subsec:profile_comparison} reports P1--P5 separately.

\subsection{Benchmark Fragmentation and the Representativeness Gap}
\label{subsec:gap}

Current benchmarks operationalise valuable local notions of competence and safety~\cite{mohammadi2025agenteval}: AgentBench, WebArena, SWE-bench cover interactive task completion and software engineering~\cite{liu2023agentbench,zhou2023webarena,jimenez2023swebench}; HaluEval and JailbreakBench probe hallucination and jailbreak robustness~\cite{li2023halueval,chao2024jailbreakbench}; API-Bank tests tool use~\cite{li2023apibank}; and GAIA, WorkArena, OSWorld, $\tau$-bench, ToolSandbox, AgentDojo extend coverage to general assistants, enterprise software, computer-use, and prompt-injection robustness~\cite{mialon2023gaia,drouin2024workarena,xie2024osworld,yao2024taubench,lu2024toolsandbox,debenedetti2024agentdojo}. But each suite targets a particular slice and a particular subset of P1--P5: they are collectively \emph{capability-isolated} (coding, planning, hallucination, safety tested separately when real failures emerge from their interaction), \emph{environment-simplified} (clean tools vs.\ real noise, latency, partial observability), and \emph{socially decontextualised} (delegation, manipulation, risk asymmetry abstracted away). An illustrative audit of representative benchmarks against eight deployment-relevant dimensions confirms strong axis-specialisation: prior suites form valuable \emph{benchmark islands} but collectively under-cover deployment-critical dimensions---operational constraints, social context, and consequence-sensitive risk---which correspond to P3--P5.

We call this the \emph{representativeness gap}. Let $\mathcal{S}$ denote the space of socio-technical scenarios relevant to deployment, $P_{\mathrm{deploy}}$ the real deployment distribution over $\mathcal{S}$, and $\mathbf{m}(a,s)=(m_{P_1}(a,s),\dots,m_{P_5}(a,s))$ a per-property trustworthiness vector for agent $a$ in scenario $s$. Deployment-level trustworthiness is the property-wise expectation $\mathbf{T}(a) = \mathbb{E}_{s \sim P_{\mathrm{deploy}}}[\mathbf{m}(a,s)]$. Benchmarks estimate from a narrower $P_{\mathrm{bench}}$ with limited long-tail coverage; when $P_{\mathrm{bench}}$ is misaligned with $P_{\mathrm{deploy}}$, scores become unreliable estimators of real trustworthiness, especially for low-frequency, high-consequence scenarios~\cite{eriksson2025trustbenchmarks}. The bottleneck is no longer metric design alone, but \emph{distribution design}~\cite{liang2023helm,kiela2021dynabench,shirali2022dynamicbench,koh2021wilds}: the missing object is a \emph{representative, risk-sensitive scenario distribution} jointly evaluated across all five properties.

\subsection{Provenance and Delta}
\label{subsec:provenance_delta}

The closest adjacent works to HAAF are two tool-use evaluation suites, AgentDojo~\cite{debenedetti2024agentdojo} and ToolSandbox~\cite{lu2024toolsandbox}. Table~\ref{tab:related_delta} contrasts all three along six deployment-relevant axes.

\begin{table*}[!htbp]
\centering
\caption{Adjacent tool-use benchmarks vs.\ HAAF. ``Props'' = deployment-level trustworthiness properties (P1--P5) explicitly scored; ``R/B Cycle'' = iterative red-team$\,\to\,$blue-team$\,\to\,$re-evaluation loop; ``IR'' = Improper Refusal as a first-class evaluated class. PI = prompt injection; fams.\ = families; dom.\ = domains; L4 = Layer-4 subset-selected.}
\label{tab:related_delta}
\small
\setlength{\tabcolsep}{3pt}
\resizebox{\linewidth}{!}{%
\begin{tabular}{@{}lcccccl@{}}
\toprule
\textbf{Benchmark} & \textbf{Props} & \textbf{Models} & \textbf{R/B} & \textbf{IR} & \textbf{Sandbox} & \textbf{Scenario Source} \\
\midrule
AgentDojo~\cite{debenedetti2024agentdojo}     & 1 (PI only)        & 10              & no  & no  & 74 tools, 4 dom.   & 97 tasks + 629 inj., hand-auth.\\
ToolSandbox~\cite{lu2024toolsandbox}          & 0 (5 cap.\ cats.)  & 13              & no  & no  & 34 tools, 11 dom.  & 1{,}032 cases, hand+seeded \\
HAAF (ours)                                   & 5 (P1--P5+IR)      & 13 (7 fams.)    & yes & yes & 5 synth.\ tools    & 100 hand-auth.\ (L4 subset)\\
\bottomrule
\end{tabular}}
\end{table*}

\paragraph{Provenance and delta.}
HAAF inherits the synthetic-sandbox paradigm of AgentDojo and ToolSandbox but extends it on three axes the adjacent suites leave out. First, \emph{trustworthiness coverage}: AgentDojo scores one property (PI robustness) and ToolSandbox scores zero (its five categories are tool-use \emph{capability} dimensions); HAAF scores all five P1--P5 jointly on the same trajectory population, and promotes Improper Refusal to a first-class evaluated class under P1. Second, \emph{model breadth}: prior suites evaluate fixed lists (10 and 13) at one time slice, whereas HAAF spans 13 models across 7 families with version-pair contrasts that expose anti-scaling regressions (\S\ref{subsec:antiscaling_stats}). Third, \emph{iterative protocol}: AgentDojo ships four static defenses and ToolSandbox has no defense loop; HAAF closes a red-team$\,\to\,$blue-team$\,\to\,$re-evaluation cycle with paired-bootstrap significance testing.
